\documentclass[11pt, a4paper]{article}
\input{preamble.tex}
\switchlang{en}
\begin{document}

\title{1980 -- A2D Joystick for Apple II, 1st generation}
\date{}
\author{~}
\maketitle
\selectlanguage{english}

The A2D Joystick (fig. \ref{fig:A2DJoystick1stGenPic}), the first-generation analog joystick from the A2D Company for the Apple II computers, was initially developed in the late 1970s and is associated with Jef Raskin, a renowned engineer and interface researcher who played a key role in the creation of the Apple Macintosh. In \cite{info_3}, when describing the second generation A2D joystick, a previous model is mentioned, referenced as the A2D 1002 model, but other mentions of this model numbers were not found.

\begin{figure}[h]
   \centering
    \includegraphics[width=\textwidth]{1980_a2d_joystick_1st_gen/pic_30.jpg}
    \caption{A2D Joystick}
    \label{fig:A2DJoystick1stGenPic}
\end{figure}

The device borrows the appearance and internal design from aircraft R/C models remote controls of the same time. Following their example, the device is enclosed into a ``heavy gold vinyl clad aluminum case (not vac formed)'' \cite{joystick81}. As seen in fig. \ref{fig:A2DJoystick1stGenTopAndBottom}, the ``heavy aluminum case'' has rectangular shape, and its underside is completely flat, without any elements, including feet. On the top side, there is a plate with the manufacturer's logo, as well as the stick of the analog joystick unit. At one end of the body there are two cut-in red tact buttons. At the other end there are two rotary potentiometer handles that allow user to set zero voltage for each of the axes (fig. \ref{fig:A2DJoystick1stGenSize}). Controls to disable the auto-centering of the stick are not provided.

Both externally and internally, the joystick reproduces a sample of the very first joystick for the Apple II, created by Jeff Raskin \cite{raskin}, and based on this the A2D joystick can be considered it's serial adaptation. The only difference which can be noticed is related to the rotary controls: they had smaller size and were randomly cut into the top cover in the Raskin's prototype, and are reasonably placed to the end of the body in the A2D Company's exemplar.

\begin{figure}[h]
    \centering
    \includegraphics[scale=0.6]{1980_a2d_joystick_1st_gen/top_30.jpg}
    \includegraphics[scale=0.6]{1980_a2d_joystick_1st_gen/bottom_30.jpg}
    \caption{A2D Joystick, top and bottom views}
    \label{fig:A2DJoystick1stGenTopAndBottom}
\end{figure}

Taking into account the date of the first known mention in literature \cite{joystick80} and the fact that ``about 500 have been in use for 1 year'' stated in \cite{joystick81}, the device appearance on the marked can be dated back to 1980.

\begin{figure}[h]
    \centering
    \includegraphics[scale=0.5]{1980_a2d_joystick_1st_gen/size_30.jpg}
    \caption{A2D Joystick on a graduated pad with a grid step of 1~cm}
    \label{fig:A2DJoystick1stGenSize}
\end{figure}

The joystick assembly, including the entire electromechanical unit, was also borrowed from R/C models remote controls. Literature emphasize this relation, mentioning ``precision open-gimbaled self-centering stick'' \cite{joystick80}, which was ``designed for radio control use'' \cite{joustickAD} and ``results in a lighter feel and more precise movement'' \cite{info_3}. In fact, the stick is quite comfortable to control with fingers (fig. \ref{fig:A2DJoystick1stGenHand}). The same cannot be said for the miniature buttons without clear tactile feedback -- ones that Jef Raskin, not yet a world authority in the world of graphical interfaces, unwisely borrowed from those same R/C remote controls. The buttons don't produce much load when controlling an airplane model, but the situation is exactly opposite for a computer joystick positioned as ``nice hand-held design with two game control push buttons'' \cite{joystick81}. Using it with computer games and other typical software any long undoubtedly creates a lot of discomfort. Apparently, this flaw was quickly noticed by the manufacturer and fixed in the next model, the A2D Joystick 2001, which received a plastic body with rounded corners and concave keycaps.

\begin{figure}[h]
    \centering
    \includegraphics[scale=0.5]{1980_a2d_joystick_1st_gen/hand_30.jpg}
    \caption{A2D Joystick with a human hand model}
    \label{fig:A2DJoystick1stGenHand}
\end{figure}

The internal design of the joystick assembly is fairly standard, based on analog potentiometers, converting the deflection of the stick into a change in the electrical resistance (fig. \ref{fig:A2DJoystick1stGenInside}). Like in the remote controls of R/C models, the sticks of A2D Joysticks are indeed placed on the metal gimbals actively advertised by the manufacturer.

\begin{figure}[h]
    \centering
    \includegraphics[scale=0.57]{1980_a2d_joystick_1st_gen/inside_30.jpg}
    \caption{A2D Joystick disassembled}
    \label{fig:A2DJoystick1stGenInside}
\end{figure}

\begin{thebibliography}{9}

\bibitem{joystick81} Van Diver G., Love R. Apple II/III software directory. Vital Information, Inc., 1981. p. PR-21 \url{https://archive.org/details/vanlovesappleiii00gera/page/470/mode/1up?q=a2d+joystick}

\bibitem{raskin} Jef Raskin's Apple II Joystick // DigiBarn Devices: Jef Raskin's Apple II Joystick \url{https://www.digibarn.com/collections/devices/raskin-joystick/index.html"}

\bibitem{joystick80}Frederick F.J. Guide to microcomputers / Association for Educational Communications and Technology. AECT-ERIC/Information Resources, Syracuse University, 1980. p. 37 \url{https://archive.org/details/guidetomicrocomp0000fred/page/37/mode/1up?q=a2d}

\bibitem {joustickAD} A2D Joystick 2001 (advertising) // Creative Computing, Vol.8, No 12, December 1982 -- p. 238 \url{https://archive.org/details/CreativeComputing1982-12/page/n231/mode/2up}

\bibitem {info_3} Ahl D.H. Joysticks, Paddles, and Game Port Extenders. Potentiometer Joysticks. A2D Joystick (2001) // Creative Computing, Vol.8, No 8, August 1982 -- p. 79. \url{https://archive.org/details/creativecomputing-1982-08/page/n79/mode/2up} % 82

\end{thebibliography}
\end{document}
